\section{Reject impossible counts before reading pointers}
The head reader accepts any unsigned available count. Zero fails with status one before
pointer access. For a nonzero count, it reads the initial byte, establishes a head size
from one through nine, and compares that size unsigned before reading argument bytes.
A large available count does not make an unsupported type valid.

The skipper first rejects an unsigned count above 1024 with status two, then rejects
inherited depth above eight with status one, before pointer access. A zero count also
fails with one. The decoder rejects counts above 1024 with two and below seven with
one before reading either input or output. State these early checks in this order so
an excessive count has a specified result even when another argument is also invalid.

For a flat record, the bound and encoder first read the sixteen-byte counter/length
prefix. Reject either length outside its field bound with status one before reading
identifier or sample storage. Every bit pattern in the counter is permitted. An invalid
length is not repaired by clipping it to the available room; doing so would encode a
different record while reporting success.
